% !PRINTED 110
\weberpage{110}

Then we first have to show that
\begin{equation}
 (y+iz)^2=b+ic,
\tag{3}
\end{equation}
whatever real values $b,c$ may have, can be satisfied by real $y,z$. Equation (3), however, is fulfilled if
\begin{equation}
 y^2-z^2=b,
\qquad 2yz=c,
\tag{4}
\end{equation}
and hence
\[
 (y^2+z^2)^2=b^2+c^2
\]
therefore
\begin{equation}
 y^2+z^2=\sqrt{b^2+c^2}.
\tag{5}
\end{equation}
By IX there always exists a positive square root $\sqrt{b^2+c^2}$. From (4) and (5) it now follows that
\[
 y^2=\frac{b+\sqrt{b^2+c^2}}{2},\qquad
 z^2=\frac{-b+\sqrt{b^2+c^2}}{2},
\]
and the two magnitudes $\pm b+\sqrt{b^2+c^2}$ are plainly positive, since $b^2+c^2>b^2$, hence also $\sqrt{b^2+c^2}>\sqrt{b^2}$.% [sic] The strict inequalities fail when c=0: then b^2+c^2=b^2 and one of the two magnitudes is 0 rather than positive. Non-negativity suffices for the construction; the source-original error is preserved rather than silently normalized.
\ Thus we obtain
\[
 y=\pm\sqrt{\frac{b+\sqrt{b^2+c^2}}{2}},\qquad
 z=\pm\sqrt{\frac{-b+\sqrt{b^2+c^2}}{2}},
\]
but one sign is determined by the other through the condition
\[
 2yz=c;
\]
one of the two signs remains arbitrary, so that we obtain not merely one, but two opposite solutions of (3).

Further let $n$ be odd; we take the equation to be solved in the form
\begin{equation}
 (y+iz)^n=b+ci.
\tag{6}
\end{equation}
If $c=0$, so that $a$ is real, then we can put $z=0$ and by IX obtain a real value of $x$. But if we assume $c$ different from zero, then from (6) (Introduction, p. 19) there follows
\begin{equation}
 (y-iz)^n=b-ci.
\tag{7}
\end{equation}
Multiplying the two equations (6), (7) by one another gives
\[
 (y^2+z^2)^n=b^2+c^2,
\]
and hence, by IX,
%CONT>>
